\documentclass[11pt,reqno]{amsart}

\usepackage[T1]{fontenc}
\usepackage{amsmath,amssymb,amsthm}
\usepackage[letterpaper,margin=1in]{geometry}
\usepackage{microtype}
\usepackage{xcolor}
\usepackage{hyperref}
\usepackage{cleveref}

\definecolor{linkblue}{RGB}{35,72,139}
\hypersetup{
  colorlinks=true,
  linkcolor=linkblue,
  citecolor=linkblue,
  urlcolor=linkblue,
  pdftitle={Short proof of the Bernoulli conjecture}
}

\AtBeginDocument{
  \setlength{\abovedisplayskip}{8pt plus 2pt minus 2pt}
  \setlength{\belowdisplayskip}{8pt plus 2pt minus 2pt}
  \setlength{\abovedisplayshortskip}{4pt plus 2pt}
  \setlength{\belowdisplayshortskip}{5pt plus 2pt minus 2pt}
}

\newtheorem{theorem}{Theorem}[section]
\newtheorem{lemma}[theorem]{Lemma}
\crefname{theorem}{Theorem}{Theorems}
\crefname{lemma}{Lemma}{Lemmas}

\newcommand{\E}{\mathbb E}
\newcommand{\Pp}{\mathbb P}
\newcommand{\R}{\mathbb R}
\newcommand{\cx}{\preceq_{\mathrm{cx}}}
\newcommand{\ip}[2]{\langle #1,#2\rangle}

\title[Short proof of the Bernoulli conjecture]
{Short proof of the Bernoulli conjecture}
\author{}
\date{August 11, 2026}
\subjclass[2020]{60G50, 60G15, 60E15, 46B09}
\keywords{Bernoulli theorem, convex comparison, Gaussian processes,
randomized rounding}

\begin{document}

\begin{abstract}
We give a short proof of the Bernoulli conjecture. The proof was produced
by an internal model at OpenAI.
\end{abstract}

\maketitle

\section{Introduction}

Let $\varepsilon=(\varepsilon_i)_{i\in I}$ have independent uniform sign
coordinates, and let $G=(g_i)_{i\in I}$ have independent standard Gaussian
coordinates. For a nonempty $T\subset\ell^2(I)$, write
\[
 b(T)=\E\sup_{t\in T}\sum_{i\in I}t_i\varepsilon_i,
 \qquad
 g(T)=\E\sup_{t\in T}\sum_{i\in I}t_i g_i,
\]
with the usual finite-subset interpretation. Recall that $X\cx Y$ means
$\E\Phi(X)\leq\E\Phi(Y)$ for every convex $\Phi$ for which the expectations
are defined.

\begin{theorem}[Bernoulli conjecture]\label{thm:bernoulli}
There is a universal constant $C$ such that, whenever
$\varnothing\ne T\subset\ell^2(I)$ and $b(T)<\infty$, there exist
$U\subset\ell^1(I)$ and $V\subset\ell^2(I)$ such that
\[
 T\subset U+V,
 \qquad
 \sup_{u\in U}\|u\|_1\leq Cb(T),
 \qquad
 g(V)\leq Cb(T).
\]
\end{theorem}

The conjecture was proposed by Talagrand~\cite{Talagrand} and first proved
by Bednorz and Lata\l a~\cite{BednorzLatala}. The main input for our proof
is a subgaussian convex-comparison theorem, which follows by duality from
Liu's comparison principle~\cite[Corollary~2]{LiuLifting} and was
explicitly formulated in convex order by van Handel~\cite{vanHandel}.
In fact, the key input for \cref{thm:bernoulli} is the following convex
domination result.

\begin{theorem}[Gaussian-to-Bernoulli comparison]\label{thm:comparison}
There is a universal constant $C_0$ such that, for every
$G\sim N(0,I_n)$ and every uniform
$\varepsilon\in\{-1,1\}^n$,
\[
 X\cx G,\qquad \|X\|_\infty\leq1\quad\text{almost surely}
 \quad\Longrightarrow\quad
 X\cx C_0\varepsilon.
\]
\end{theorem}

We use this Gaussian comparison in the following form: for a
universal $C_G$, if $Y$ is centered and
\begin{equation}\label{eq:subgaussian}
 \log\E e^{\ip{a}{Y}}
 \leq\frac{K^2}{2}\|a\|_2^2
 \qquad(a\in\R^n),
\end{equation}
then $Y\cx C_GK G'$, where $G'\sim N(0,I_n)$
\cite{vanHandel}. By Strassen's theorem~\cite{Strassen}, this is
equivalent to the existence of a martingale coupling with
\begin{equation}\label{eq:gaussian-martingale}
 Y=C_GK\E[G'\mid Y].
\end{equation}

\smallskip
\noindent\textbf{Acknowledgment.}
We thank Ramon van Handel for various helpful comments.

\smallskip
\noindent\textbf{Note on concurrent work:}
During the preparation of this manuscript, we learned of an independent
alternate proof of the Bernoulli conjecture by Liu and
Zadik~\cite{LiuZadik}.

\section{Reduction to
  \texorpdfstring{\cref{thm:comparison}}
                 {Theorem \getrefnumber{thm:comparison}}}
\label{sec:reduction}

We first give the reduction from \cref{thm:bernoulli} to
\cref{thm:comparison}.

\begin{proof}
By the standard finite-section and compactness reduction, it suffices to
prove \cref{thm:bernoulli} for a finite
$T\subset\R^n$. Set
\begin{equation}\label{eq:D}
 D(T)=\inf_{(v_t)_{t\in T}}
 \left\{
   \max_{t\in T}\|t-v_t\|_1
   +\E\max_{t\in T}\ip{v_t}{G}
 \right\}.
\end{equation}
For a measurable probability kernel $q=(q_t)_{t\in T}$, put
\[
 q_t\geq0,\qquad \sum_{t\in T}q_t=1,
 \qquad
 m_t=\E[Gq_t(G)].
\]
We claim that
\begin{equation}\label{eq:minimax}
 D(T)=
 \sup_{q}
 \left\{
   \sum_{t\in T}\ip{t}{m_t}:
   \sum_{t\in T}\|m_t\|_\infty\leq1
 \right\}.
\end{equation}
Indeed,
\begin{align*}
 \max_t\|t-v_t\|_1
  &=\sup_{\sum_t\|z_t\|_\infty\leq1}
      \sum_t\ip{t-v_t}{z_t},\\
 \E\max_t\ip{v_t}{G}
  &=\sup_q\sum_t\ip{v_t}{m_t}.
\end{align*}
The kernels form a weakly compact convex subset of
$L^2(\gamma_n)^{T}$, so one application of minimax gives
\[
 D(T)=
 \sup_{\substack{\sum_t\|z_t\|_\infty\leq1\\q_t\geq0,\ \sum_tq_t=1}}
 \inf_{(v_t)}
 \left\{
   \sum_t\ip{t}{z_t}
   +\sum_t\ip{v_t}{m_t-z_t}
 \right\}.
\]
The inner infimum is finite only when $z_t=m_t$ for every $t$:
if $m_s\ne z_s$, take
$v_s=-R(m_s-z_s)$ and $v_t=0$ for $t\ne s$ and let
$R\to\infty$. In the remaining case the inner expression is
$\sum_t\ip{t}{m_t}$. This proves~\eqref{eq:minimax} and, in
particular, explains precisely why the dual variable is forced to equal
the Gaussian barycenter.

Fix an admissible kernel in~\eqref{eq:minimax}, and write
$p_t=\E q_t(G)$ and $a_t=\|m_t\|_\infty$.
Since $\sum_t(a_t-p_t)_+\leq\sum_ta_t\leq1$, choose a probability
vector $r=(r_t)_{t\in T}$ such that
$r_t\geq(a_t-p_t)_+$. Toss an independent fair coin. On heads,
sample a label $J$ from $q(G)$; on tails, sample $J$ independently
of $G$ with distribution $r$. Then
\[
 \Pp(J=t)=\frac{p_t+r_t}{2},
 \qquad
 \E\bigl[G\mathbf1_{\{J=t\}}\bigr]=\frac{m_t}{2}.
\]
Consequently $X=\E[G\mid J]$ satisfies $X\cx G$, and on every
positive-probability event $\{J=t\}$,
\[
 X=\frac{m_t}{p_t+r_t},
 \qquad
 \|X\|_\infty=\frac{a_t}{p_t+r_t}\leq1.
\]
Apply \cref{thm:comparison} to the convex function
$\Phi(x)=\max_{t\in T}\ip{t}{x}$. We obtain
\[
 \frac12\sum_{t\in T}\ip{t}{m_t}
   =\E\ip{J}{X}
   \leq\E\Phi(X)
   \leq C_0b(T).
\]
Taking the supremum in~\eqref{eq:minimax} yields
$D(T)\leq2C_0b(T)$. The infimum in~\eqref{eq:D} is attained by
coercivity. For a minimizer $(v_t)$ set
$U=\{t-v_t:t\in T\}$ and $V=\{v_t:t\in T\}$.
Both terms in~\eqref{eq:D} are nonnegative. Hence
\[
 T\subset U+V,
 \qquad
 \sup_{u\in U}\|u\|_1\leq2C_0b(T),
 \qquad
 g(V)\leq2C_0b(T),
\]
as required.
\end{proof}

\section{Proof of
  \texorpdfstring{\cref{thm:comparison}}
                 {Theorem \getrefnumber{thm:comparison}}}
\label{sec:lifting}

The first lemma is the finite-support Gaussian case of a
maximum-entropy martingale bridge. It is a special case of the
construction of Hua, Song, and
Tudose~\cite[Proposition~3.3 and Section~4.1]{HuaSongTudose}, applied
to $X/2\cx G/2$. The general theory of martingale Schr\"odinger bridges
is developed by Nutz and Wiesel~\cite{NutzWiesel} and by Backhoff,
Beiglb\"ock, Bifronte, and Ley~\cite{BackhoffEtAl}. We include the
short argument for completeness.

\begin{lemma}\label[lemma]{lem:lifting}
Suppose $X$ has finite support and $X\cx G$. There exists a coupling
of $X$ and $G\sim N(0,I_n)$ such that
\begin{equation}\label{eq:half-lifting}
 \E[G\mid X]=\tfrac12X,
\end{equation}
and each conditional distribution $G\mid\{X=x\}$ has an everywhere
positive, $1$-strongly log-concave density. Among couplings satisfying
\eqref{eq:half-lifting}, this coupling uniquely maximizes $H(X\mid G)$.
\end{lemma}

\begin{proof}
Write $\Pp(X=x_j)=p_j>0$ for $1\leq j\leq m$.
Convex comparison gives $\E X=0$. Set $a_m=0$, $b_m=0$, and
$\ell_j(g)=a_j+\ip{b_j}{g}$. Consider
\begin{equation}\label{eq:partition}
 P(a,b)=
 \E\log\sum_{j=1}^m e^{\ell_j(G)}
 -\sum_{j=1}^mp_j
    \left(a_j+\tfrac12\ip{b_j}{x_j}\right).
\end{equation}
Let $h=\E\max_j\ell_j(G)$. Since log-sum-exp dominates the maximum,
\begin{equation}\label{eq:coercivity}
 P(a,b)\geq
 \tfrac12\left(h-\sum_jp_j\ell_j(x_j)\right)
 +\tfrac12\E\left[
     \max_j\ell_j(G)-\sum_jp_j\ell_j(G)
   \right].
\end{equation}
The first term is nonnegative by $X\cx G$. The second is continuous,
positively homogeneous, and strictly positive unless all $\ell_j$
coincide, hence unless $(a,b)=0$. Compactness of the unit sphere now
shows that $P$ is coercive, so it has a minimizer $(a^*,b^*)$.

Define
\begin{equation}\label{eq:gibbs}
 q_j(g)=
 \frac{e^{a_j^*+\ip{b_j^*}{g}}}
      {\sum_{k=1}^me^{a_k^*+\ip{b_k^*}{g}}}.
\end{equation}
Differentiating~\eqref{eq:partition} gives
\[
 \E q_j(G)=p_j,
 \qquad
 \E[Gq_j(G)]=\tfrac12p_jx_j
 \qquad(1\leq j<m).
\]
The identities for $j=m$ follow by summation from
$\E G=\E X=0$. Conditional on $G=g$, assign $X=x_j$ with
probability $q_j(g)$. This has the prescribed $X$ marginal and
satisfies~\eqref{eq:half-lifting}.

To see the maximum-entropy claim, let $\tilde q$ be any other kernel
with the same marginal and barycenter constraints, and write
$Z(g)=\sum_ke^{a_k^*+\ip{b_k^*}{g}}$ and
$H(\tilde q)=-\E\sum_j\tilde q_j(G)\log\tilde q_j(G)$.
Since $\log q_j(g)=a_j^*+\ip{b_j^*}{g}-\log Z(g)$, those
constraints give
\[
 \E\sum_j\tilde q_j(G)\log q_j(G)
 =\sum_jp_j\left(a_j^*+\tfrac12\ip{b_j^*}{x_j}\right)
   -\E\log Z(G)
 =\E\sum_jq_j(G)\log q_j(G).
\]
Thus the cross-entropy term is the same for every feasible kernel, and
\[
 H(q)-H(\tilde q)
 =\E\sum_j\tilde q_j(G)
       \log\frac{\tilde q_j(G)}{q_j(G)}\geq0.
\]
Equality holds only when $\tilde q=q$ almost surely. Finally,
$\log q_j$ is concave because log-sum-exp is convex. Hence the
conditional Gaussian density
$g\mapsto[p_j(2\pi)^{n/2}]^{-1}q_j(g)e^{-\|g\|_2^2/2}$
is everywhere positive and $1$-strongly log-concave.
\end{proof}

The following is an immediate consequence of the Brascamp--Lieb
variance inequality~\cite[Theorem~4.1]{BrascampLieb}.

\begin{lemma}\label[lemma]{lem:variance}
For the coupling in \cref{lem:lifting},
\[
 \operatorname{Var}(G_i\mid X)\leq1,
 \qquad
 \E[G_i^2\mid X]\leq1+\tfrac14X_i^2.
\]
\end{lemma}

We next record a standard subgaussian estimate needed for
\cref{lem:rounding}.

\begin{lemma}\label[lemma]{lem:bernstein}
Let $V_1,\ldots,V_n$ be independent and centered, with
$\E e^{|V_i|/\kappa}\leq2$, and let
$Y_i=\E[V_i\mid\mathcal H]$ satisfy $|Y_i|\leq D$. Then
\[
 \log\E e^{\ip{a}{Y}}
 \leq\max\{4\kappa^2,2\kappa D\}\|a\|_2^2
 \qquad(a\in\R^n).
\]
\end{lemma}

\begin{proof}
The standard Bernstein moment estimate for a centered subexponential
variable~\cite[proof of Proposition~2.8.1]{Vershynin} gives
\[
 \E e^{sV_i}\leq e^{4\kappa^2s^2}
 \qquad\text{if }\kappa|s|\leq\tfrac12.
\]
Set $A=\{i:\kappa|a_i|>1/2\}$. For $i\in A$,
$D|a_i|\leq2\kappa D a_i^2$. Thus conditional Jensen and independence
give
\[
 \E e^{\ip{a}{Y}}
 \leq
 e^{2\kappa D\sum_{i\in A}a_i^2}
 \prod_{i\notin A}\E e^{a_iV_i}
 \leq
 \exp\left(
  2\kappa D\sum_{i\in A}a_i^2
  +4\kappa^2\sum_{i\notin A}a_i^2
 \right).
\]
The claimed bound follows.
\end{proof}

\begin{lemma}\label[lemma]{lem:rounding}
Suppose $\|X\|_\infty\leq1$, and use the coupling of
\cref{lem:lifting}. For every $L>0$, there exist a uniform
$\varepsilon\in\{-1,1\}^n$ and a centered $Y\in\R^n$ such that
\begin{equation}\label{eq:rounding}
 \tfrac12X=L\E[\varepsilon\mid X]+Y,
 \qquad
 \|Y\|_\infty\leq\frac1{2L},
 \qquad
 \log\E e^{\ip{a}{Y}}
 \leq\frac4{L^2}\|a\|_2^2.
\end{equation}
\end{lemma}

\begin{proof}
Independently of $(G,X)$, take independent
$U_i\sim\operatorname{Unif}[-L,L]$ and set
$\varepsilon_i=\operatorname{sgn}(G_i+U_i)$.
The pairs $(G_i,U_i)$ are independent and symmetric. Therefore the
$\varepsilon_i$ are genuinely independent uniform signs; no
independence conditional on $X$ is asserted or needed. Moreover,
\[
 \E[\varepsilon_i\mid G,X]
 =\operatorname{clip}(G_i/L,-1,1).
\]
Define $V_i=\operatorname{sgn}(G_i)(|G_i|-L)_+$ and
$Y_i=\E[V_i\mid X]$.
The $V_i$ are independent and centered, and
$G_i=L\operatorname{clip}(G_i/L,-1,1)+V_i$.
Taking conditional expectations and
using~\eqref{eq:half-lifting} proves the first identity
in~\eqref{eq:rounding}.

The elementary inequality $(r-L)_+\leq r^2/(4L)$ gives
\[
 \E e^{L|V_i|}
 \leq\E e^{G_i^2/4}=\sqrt2\leq2.
\]
\cref{lem:variance} and $|X_i|\leq1$ further give
\[
 |Y_i|
 \leq\frac{\E[G_i^2\mid X]}{4L}
 \leq\frac{1+X_i^2/4}{4L}
 \leq\frac1{2L}.
\]
Apply \cref{lem:bernstein} with
$\kappa=L^{-1}$ and $D=(2L)^{-1}$ to obtain the final
bound in~\eqref{eq:rounding}.
\end{proof}

\begin{proof}[Proof of \cref{thm:comparison}]
For a convex $\phi:\R^n\to\R$, write
\[
 S_\phi=\sup\bigl\{\E\phi(X):X\cx G,\ \|X\|_\infty\leq1\bigr\}<\infty.
\]
It suffices to show that $S_\phi\leq\E\phi(C_0\varepsilon)$ for every
convex $\phi$. Conditional expectation on increasingly fine finite
partitions of the cube shows that the supremum defining $S_\phi$ may be
taken over finite-valued $X$.

Set $C_* =\max\{2\sqrt2C_G,1/2\}$, fix $L>2C_*$, and put
$\lambda=C_*/L$ and $C_0=2L/(1-2\lambda)$. For each finite-valued
admissible $X$, \cref{lem:rounding} and Gaussian
comparison~\eqref{eq:gaussian-martingale} yield a uniform sign vector
$\varepsilon$ and $G'\sim N(0,I_n)$ such that
\[
 \tfrac12X=L\E[\varepsilon\mid X]+\lambda Z,
 \qquad Z=\E[G'\mid X].
\]
Jensen gives $Z\cx G$, and \cref{lem:rounding} gives
\[
 \|Z\|_\infty\leq\frac1{2L\lambda}
 =\frac1{2C_*}\leq1,
 \qquad
 X=(1-2\lambda)\E[C_0\varepsilon\mid X]+2\lambda Z.
\]
Thus another application of Jensen gives
\[
 \E\phi(X)
 \leq(1-2\lambda)\E\phi(C_0\varepsilon)
      +2\lambda\E\phi(Z)
 \leq(1-2\lambda)\E\phi(C_0\varepsilon)+2\lambda S_\phi.
\]
Taking the supremum over $X$ and using $2\lambda<1$ yields
$S_\phi\leq\E\phi(C_0\varepsilon)$, as claimed.
\end{proof}

\begin{thebibliography}{99}

\bibitem{BackhoffEtAl}
J.~Backhoff, M.~Beiglb\"ock, G.~Bifronte, and A.~Ley,
\emph{Bridging classical and martingale Schr\"odinger bridges},
2026, \href{https://arxiv.org/abs/2604.01299}{arXiv:2604.01299}.

\bibitem{BednorzLatala}
W.~Bednorz and R.~Lata\l a,
\emph{On the boundedness of Bernoulli processes},
Ann. of Math. (2) \textbf{180} (2014), 1167--1203,
\href{https://doi.org/10.4007/annals.2014.180.3.8}
{doi:10.4007/annals.2014.180.3.8}.

\bibitem{BrascampLieb}
H.~J.~Brascamp and E.~H.~Lieb,
\emph{On extensions of the Brunn--Minkowski and
Pr\'ekopa--Leindler theorems, including inequalities for log concave
functions, and with an application to the diffusion equation},
J. Funct. Anal. \textbf{22} (1976), 366--389,
\href{https://doi.org/10.1016/0022-1236(76)90004-5}
{doi:10.1016/0022-1236(76)90004-5}.

\bibitem{HuaSongTudose}
D.~M.~Hua, A.~Song, and S.~Tudose,
\emph{On Talagrand's Convexity Conjecture},
2026, \href{https://arxiv.org/abs/2605.10908}{arXiv:2605.10908}.

\bibitem{LiuLifting}
J.~Liu,
\emph{Simple and Sharp Generalization Bounds via Lifting},
2025, \href{https://arxiv.org/abs/2508.18682}{arXiv:2508.18682}.

\bibitem{LiuZadik}
J.~Liu and I.~Zadik,
\emph{A Bayesian Proof of the Bernoulli Theorem},
2026, \href{https://arxiv.org/abs/2608.11031}{arXiv:2608.11031}.

\bibitem{NutzWiesel}
M.~Nutz and J.~Wiesel,
\emph{On the martingale Schr\"odinger bridge between two distributions},
to appear in Bernoulli,
\href{https://arxiv.org/abs/2401.05209}{arXiv:2401.05209}.

\bibitem{Strassen}
V.~Strassen,
\emph{The existence of probability measures with given marginals},
Ann. Math. Statist. \textbf{36} (1965), 423--439,
\href{https://doi.org/10.1214/aoms/1177700153}
{doi:10.1214/aoms/1177700153}.

\bibitem{Talagrand}
M.~Talagrand,
\emph{Constructions of majorizing measures, Bernoulli processes and
cotype},
Geom. Funct. Anal. \textbf{4} (1994), 660--717,
\href{https://doi.org/10.1007/BF01896658}{doi:10.1007/BF01896658}.

\bibitem{vanHandel}
R.~van Handel,
\emph{On the subgaussian comparison theorem},
2025, \href{https://arxiv.org/abs/2512.18588}{arXiv:2512.18588}.

\bibitem{Vershynin}
R.~Vershynin,
\emph{High-Dimensional Probability: An Introduction with Applications
in Data Science},
second ed., Cambridge University Press, Cambridge, 2026,
\href{https://doi.org/10.1017/9781009490672}
{doi:10.1017/9781009490672}.

\end{thebibliography}

\end{document}
